\documentclass[../../main.tex]{subfiles}
\begin{document}

\begin{center}
  {\fontsize{48pt}{58pt}\selectfont \textbf{1993}\normalsize\textbf{年度}}
\end{center}

\begin{table}[htb]
  \centering
  \caption{出題テーマと難易度}
  \begin{tabular}{|c|p{6.2cm}|c|c|c|}
    \hline
    \textbf{問題番号} & \textbf{テーマ} & \textbf{思考量} & \textbf{計算量} & \textbf{総合} \\
    \hline
    第1問 & 双曲線への接線で作る三角形の面積 & 普 & 普 & 普 \\
    \hline
    第2問 & 円内の2円への接線の長さと最大値 & 普 & 普 & 普 \\
    \hline
    第3問 & 多項式の積分方程式 & 普 & 普 & 普 \\
    \hline
    第4問 & 定積分で定めた数列の極限 & 普 & 普 & 普 \\
    \hline
    第5問 & 二項分布と確率の不等式 & 普 & 普 & 普 \\
    \hline
    第6問 & 行列による一次変換 &  &  &  \\
    \hline
  \end{tabular}
\end{table}

\begin{table}[htb]
  \centering
  \caption{解答}
  \renewcommand{\arraystretch}{1.4}
  \begin{tabular}{|c|c|p{9cm}|}
    \hline
    \textbf{問題番号} & \textbf{小問} & \textbf{答え} \\
    \hline
    第1問 & -- & $p=\sqrt{2}$ \\
    \hline
    第2問 & (1) & $PQ=\sqrt{2(1-a)(1+\cos\theta)}$ \\
    \cline{2-3}
    & (2) & $\max(PQ+PR)=2\sqrt{2-a-b}$ \\
    \hline
    第3問 & -- & $f(x)$は定数（証明問題） \\
    \hline
    第4問 & (1) & $\lim_{n\to\infty}b_n=0$ \\
    \cline{2-3}
    & (2) & $\lim_{n\to\infty}na_n=1$ \\
    \hline
    第5問 & -- & $p_n=\dfrac{1}{2}\left(1+\dfrac{{}_{2n}C_n}{4^n}\right)\ge\dfrac{1}{2}+\dfrac{1}{4n}$（証明問題） \\
    \hline
    第6問 & (1) &  \\
    \cline{2-3}
    & (2) &  \\
    \cline{2-3}
    & (3) &  \\
    \hline
  \end{tabular}
\end{table}

\paragraph{試験形態}
6問．試験時間は要確認．

\end{document}
